\documentclass[12pt,a4paper]{article}
\usepackage{goyabase}

\usepackage{longtable}
\usepackage{calc}

% Pandoc compatibility
\providecommand{\tightlist}{%
  \setlength{\itemsep}{0pt}\setlength{\parskip}{0pt}}
\newcommand{\passthrough}[1]{#1}

\title{\textbf{Cassandra vs SQL: El cambio de mentalidad (Query-First)}}
\author{Departamento de Informática}
\date{\today}

\usepackage[spanish, es-noshorthands, es-tabla]{babel}
\usepackage[autostyle]{csquotes}
\begin{document}

\maketitle

\section{Cassandra vs SQL: El cambio de mentalidad (Query-First)}\label{cassandra-vs-sql-el-cambio-de-mentalidad-query-first}

\begin{center}\rule{0.5\linewidth}{0.5pt}\end{center}

\subsection{Contexto de Negocio: Registro de Búsquedas (GloboTour)}\label{contexto-de-negocio-registro-de-buxfasquedas-globotour}

En la plataforma \textbf{GloboTour}, miles de usuarios realizan búsquedas de destinos cada segundo. Guardar esto en una base de datos relacional (como MySQL) bloquearía las tablas debido al alto volumen de escritura.

Usamos \textbf{Cassandra} para este módulo de \enquote{Logs de Búsqueda} porque:
\begin{enumerate}
    \item \textbf{Escalabilidad}: Podemos añadir servidores fácilmente si el tráfico sube.
    \item \textbf{Escritura ultra-rápida}: Cassandra está diseñada para insertar millones de registros sin despeinarse.
    \item \textbf{Consulta específica}: El equipo de Marketing necesita ver rápidamente el historial de un usuario para ofrecerle ofertas personalizadas.
\end{enumerate}

\subsubsection{El Modelo de Datos Creado}

La tabla \texttt{user\_search\_logs} está diseñada para responder a: \textbf{\enquote{¿Qué ha buscado el usuario X recientemente?}}

\begin{itemize}
    \item \textbf{Partition Key (\texttt{user\_id})}: Agrupa todas las búsquedas de un usuario en el mismo nodo.
    \item \textbf{Clustering Key (\texttt{search\_time})}: Ordena esas búsquedas automáticamente de la más reciente a la más antigua.
\end{itemize}

\begin{center}\rule{0.5\linewidth}{0.5pt}\end{center}

\subsection{1. La Gran Diferencia: ¿Dónde están mis JOINs?}\label{la-gran-diferencia-duxf3nde-estuxe1n-mis-joins}

En SQL normalizamos para evitar duplicar datos. En Cassandra \textbf{la duplicación es tu amiga} para ganar velocidad.

\subsubsection{El problema del JOIN}

\begin{itemize}
    \item \textbf{SQL}: Tienes una tabla \texttt{usuarios} y una tabla \texttt{busquedas}. Para saber qué búsquedas hizo \enquote{Víctor}, haces un \texttt{JOIN}.
    \item \textbf{Cassandra}: \textbf{NO existen los JOINs}. Cruzar datos entre nodos en un cluster de 100 servidores sería lentísimo.
\end{itemize}

\subsubsection{La solución: Desnormalización (Duplicar info)}

Si quieres mostrar el nombre del usuario junto a su búsqueda, \textbf{guardas el nombre del usuario dentro de la tabla de búsquedas}, aunque se repita 1000 veces.

\begin{center}\rule{0.5\linewidth}{0.5pt}\end{center}

\subsection{Comparativa de Modelado (Caso GloboTour)}\label{comparativa-de-modelado-caso-globotour}

Imagina que queremos saber: 1. \enquote{Todas las búsquedas del usuario X}. 2. \enquote{Todas las búsquedas que se han hecho al destino \enquote{París}}.

\subsubsection{En SQL (Relacional)}

Tendrías una sola tabla y harías:

\begin{lstlisting}[language=SQL]
SELECT * FROM busquedas WHERE usuario_id = 'user1'; -- Indice en usuario_id
SELECT * FROM busquedas WHERE destino = 'Paris';   -- Indice en destino
\end{lstlisting}

\subsubsection{En Cassandra (Columnar)}

Necesitas \textbf{DOS TABLAS} para los mismos datos:

\paragraph{Tabla A: Optimizada para buscar por USUARIO}

\begin{lstlisting}[language=CQL]
CREATE TABLE busquedas_por_usuario (
    usuario_id text,
    fecha timestamp,
    destino text,
    nombre_usuario text,  -- DATO DUPLICADO (Desnormalizado)
    PRIMARY KEY (usuario_id, fecha)
) WITH CLUSTERING ORDER BY (fecha DESC);
\end{lstlisting}

\paragraph{Tabla B: Optimizada para buscar por DESTINO}

\begin{lstlisting}[language=CQL]
CREATE TABLE busquedas_por_destino (
    destino text,
    fecha timestamp,
    usuario_id text,
    PRIMARY KEY (destino, fecha)
);
\end{lstlisting}

\begin{center}\rule{0.5\linewidth}{0.5pt}\end{center}

\subsection{¿Qué pasa si necesito un JOIN \enquote{sí o sí}?}\label{quuxe9-pasa-si-necesito-un-join-suxed-o-suxed}

Si te encuentras en una situación donde necesitas datos de dos tablas, tienes tres caminos en Cassandra:

\begin{enumerate}
    \item \textbf{Duplicar la información (Recomendado)}: Añade las columnas que necesites de la Tabla B en la Tabla A en el momento de la inserción. (Espacio en disco barato vs Tiempo de CPU caro).
    \item \textbf{JOIN en la Aplicación}: Tu programa (Python/Java) hace una consulta a la Tabla A, obtiene un ID, y luego hace otra consulta a la Tabla B. (Solo para datos que cambian poco).
    \item \textbf{Tablas Complementarias}: Crear una tabla específica para esa \enquote{vista} concreta que necesitas.
\end{enumerate}

\begin{center}\rule{0.5\linewidth}{0.5pt}\end{center}

\subsection{¿Hay Foreign Keys (FK) en Cassandra?}\label{hay-foreign-keys-fk-en-cassandra}

\textbf{NO.} No existe la integridad referencial.

\begin{itemize}
    \item \textbf{¿Por qué?}: En un sistema distribuido, comprobar si un ID existe en otro nodo antes de insertar un dato sería extremadamente lento.
    \item \textbf{Consecuencia}: Puedes insertar una búsqueda para un \texttt{usuario\_id} que no existe.
    \item \textbf{Responsabilidad}: La integridad de los datos recae ahora en el \textbf{programador} (la lógica de la aplicación), no en la base de datos.
\end{itemize}

\begin{center}\rule{0.5\linewidth}{0.5pt}\end{center}

\subsection{Ejemplos de consultas que NO puedes hacer (y por qué)}\label{ejemplos-de-consultas-que-no-puedes-hacer-y-por-quuxe9}

\subsubsection{Error 1: Filtrar por algo que no es Clave de Partición}

\begin{lstlisting}[language=CQL]
-- ESTO DARA ERROR
SELECT * FROM busquedas_por_usuario WHERE destino = 'Madrid';
\end{lstlisting}

\begin{itemize}
    \item \textbf{Por qué}: Cassandra no sabe en qué nodo está \enquote{Madrid}. Tendría que escanear todo el cluster.
    \item \textbf{Solución}: Crear la tabla \texttt{busquedas\_por\_destino}.
\end{itemize}

\subsubsection{Error 2: Usar el Operador \enquote{OR}}

\begin{lstlisting}[language=CQL]
-- ESTO NO EXISTE EN CASSANDRA
SELECT * FROM busquedas_por_usuario WHERE usuario_id = 'u1' OR usuario_id = 'u2';
\end{lstlisting}

\begin{itemize}
    \item \textbf{Por qué}: Cada usuario puede estar en un servidor distinto. Las consultas deben ser precisas.
    \item \textbf{Solución}: Usar el operador \texttt{IN} o hacer dos consultas separadas.
\end{itemize}

\subsubsection{Error 3: El \enquote{SELECT COUNT(*)}}

\begin{lstlisting}[language=CQL]
-- ESTO ES MUY LENTO
SELECT COUNT(*) FROM busquedas_por_usuario;
\end{lstlisting}

\begin{itemize}
    \item \textbf{Por qué}: Cassandra tiene que recorrer todos los nodos.
    \item \textbf{Alternativa (Contadores)}: Crear una tabla específica de tipo \texttt{counter}.
\end{itemize}

\begin{lstlisting}[language=CQL]
CREATE TABLE estadisticas_destino (
    destino text PRIMARY KEY,
    total counter
);

-- Actualizar el contador (no se usa INSERT, sino UPDATE)
UPDATE estadisticas_destino SET total = total + 1 WHERE destino = 'Paris';
\end{lstlisting}

\begin{center}\rule{0.5\linewidth}{0.5pt}\end{center}

\subsection{Consultas que SÍ puedes hacer (y cómo pensar en ellas)}\label{consultas-que-si-puedes-hacer}

\begin{tcolorbox}[title=Paso previo obligatorio, colframe=red!75!black]
Antes de lanzar cualquier consulta en la consola (\texttt{cqlsh}), debes seleccionar el Keyspace de la empresa:
\begin{lstlisting}[language=CQL]
USE globotour;
\end{lstlisting}
Si no lo haces, Cassandra te devolverá el error: \textit{\enquote{No keyspace has been specified}}.
\end{tcolorbox}

Para que una consulta funcione en Cassandra, siempre debemos empezar por la \textbf{Partition Key}. En nuestra tabla \texttt{busquedas\_por\_usuario}, la PK es \texttt{user\_id}.

\subsubsection{Buscar todas las búsquedas de un usuario específico}
Esta es la consulta para la que se diseñó la tabla. Es instantánea.

\begin{lstlisting}[language=CQL]
SELECT * FROM busquedas_por_usuario 
WHERE user_id = 'user_1';
\end{lstlisting}

\subsubsection{Buscar en un rango de tiempo (Clustering Column)}
Como \texttt{search\_time} es una clustering column, podemos hacer filtros de rango \textbf{siempre que hayamos fijado el usuario primero}.

\begin{lstlisting}[language=CQL]
-- Búsquedas de user_1 realizadas después del 15 de abril
SELECT * FROM busquedas_por_usuario 
WHERE user_id = 'user_1' 
  AND search_time > '2026-04-15 00:00:00';
\end{lstlisting}

\subsubsection{Limitar resultados (Top N)}
Gracias al orden físico (\texttt{CLUSTERING ORDER BY search\_time DESC}), obtener las últimas búsquedas es muy eficiente.

\begin{lstlisting}[language=CQL]
-- Las 3 últimas búsquedas de user_7
SELECT * FROM busquedas_por_usuario 
WHERE user_id = 'user_7' 
LIMIT 3;
\end{lstlisting}

\subsubsection{El truco del ALLOW FILTERING (Uso con precaución)}
Si realmente necesitas buscar por algo que no es clave (como el destino), Cassandra te permite forzarlo, pero te advierte que será lento en un entorno real.

\begin{lstlisting}[language=CQL]
-- Buscar búsquedas a Madrid de CUALQUIER usuario
SELECT * FROM busquedas_por_usuario 
WHERE destino = 'Madrid' 
ALLOW FILTERING;
\end{lstlisting}

\begin{center}\rule{0.5\linewidth}{0.5pt}\end{center}

\subsection{Visualización Profesional con CloudBeaver}\label{visualizacion}
Para ver cómo se organizan los datos realmente en el clúster:
\begin{enumerate}
    \item Abre \url{http://localhost:8978}.
    \item Selecciona la conexión \textbf{GloboTour}.
    \item Explora la tabla \texttt{busquedas\_por\_usuario}. Verás que los datos están agrupados por partición, lo cual es la clave del rendimiento de Cassandra.
\end{enumerate}

\begin{center}\rule{0.5\linewidth}{0.5pt}\end{center}

\subsection{Desafíos de Aprendizaje (Para practicar en clase)}\label{desafuxedos-de-aprendizaje-para-practicar-en-clase}

\begin{enumerate}
    \item \textbf{La Consulta Prohibida}: Intenta buscar una búsqueda filtrando solo por \texttt{destino} en la tabla \texttt{busquedas\_por\_usuario}. ¿Qué error te da Cassandra? ¿Cómo lo solucionarías profesionalmente?
    \item \textbf{Orden Físico}: Inserta tres búsquedas para el mismo usuario pero con fechas muy distintas. Haz un \texttt{SELECT}. ¿En qué orden aparecen? ¿Por qué?
    \item \textbf{Integridad entre Tablas}: Si duplicamos los datos en dos tablas (\texttt{busquedas\_por\_usuario} y \texttt{busquedas\_por\_destino}), ¿cómo asegura Cassandra que si borro un registro en una se borre en la otra?
    \item \textbf{Ampliación Avanzada}: Una vez dominados estos conceptos básicos, dirígete al \textbf{02-Laboratorio\_QueryFirst.tex} para aprender patrones industriales como Series Temporales y TTL.
\end{enumerate}

\begin{center}\rule{0.5\linewidth}{0.5pt}\end{center}

\subsection{Solucionario de Desafíos}\label{solucionario-desafios}

\begin{tcolorbox}[title=Solución Desafío 1: El Filtro Imposible]
Si ejecutas \texttt{SELECT * FROM busquedas\_por\_usuario WHERE destino = 'Madrid';}, obtendrás un error indicando que Cassandra no puede realizar el escaneo (\textit{Cannot execute this query as it might involve data filtering and thus may have unpredictable performance}).

\textbf{Solución Profesional:} Crear una tabla específica para esa pregunta:
\begin{lstlisting}[language=CQL]
CREATE TABLE busquedas_por_destino (
    destino text,
    search_time timestamp,
    user_id text,
    PRIMARY KEY (destino, search_time)
);
\end{lstlisting}
\end{tcolorbox}

\begin{tcolorbox}[title=Solución Desafío 2: El Orden de los Datos]
Al hacer el \texttt{SELECT}, las búsquedas aparecerán ordenadas de \textbf{más reciente a más antigua}.

\textbf{Por qué:} En la definición de la tabla incluimos la cláusula \texttt{WITH CLUSTERING ORDER BY (search\_time DESC)}. Cassandra guarda los datos físicamente en ese orden dentro del disco de cada nodo, lo que hace que recuperar el historial reciente sea extremadamente rápido.
\end{tcolorbox}

\begin{tcolorbox}[title=Solución Desafío 3: Coherencia y ACID]
\textbf{Respuesta Corta:} No lo hace. Cassandra no tiene claves ajenas (FK) ni transacciones ACID entre tablas.

\textbf{Explicación:} Si borras un dato en una tabla, la otra seguirá teniéndolo a menos que tu aplicación haga los dos borrados. Para asegurar que ambos cambios ocurren \enquote{a la vez} o no ocurre ninguno, Cassandra ofrece los \textbf{BATCHES} (lotes), aunque su uso excesivo puede bajar el rendimiento. En el mundo NoSQL, se suele aceptar la \enquote{Consistencia Eventual}. ¿Recuerdas lo que es?
\end{tcolorbox}

\end{document}
